%% =========================================================
%% PRX Supplemental Material Template
%% =========================================================

\documentclass[aps,prx,preprint,superscriptaddress,nofootinbib]{revtex4-2}

%% Packages (Keep consistent with main text)
\usepackage{graphicx}
\usepackage{bm}
\usepackage{amsmath}
\usepackage{amssymb}
\usepackage{physics}
\usepackage{hyperref}
\hypersetup{colorlinks=true, linkcolor=blue, citecolor=blue, urlcolor=blue}

\begin{document}

%% =========================================================
%% Title Area
%% =========================================================
% It's common to repeat the title with "Supplemental Material for"

\title{Supplemental Material:\\Transient Freezing Governs Flat-Minimum Selection in Stochastic Gradient Descent}

\author{Ning Yang}
\thanks{These authors contributed equally to this work.}
\affiliation{Peking University Chengdu Academy for Advanced Interdisciplinary Biotechnologies, Chengdu 610213, China}

\author{Yikuan Zhang}
\thanks{These authors contributed equally to this work.}
\affiliation{School of Physics, Peking University, Beijing 100871, China}

\author{Qi Ouyang}
\affiliation{Institute for Advanced Study in Physics, Zhejiang University, Hangzhou 310058, China}

\author{Chao Tang}
\affiliation{Peking University Chengdu Academy for Advanced Interdisciplinary Biotechnologies, Chengdu 610213, China}
\affiliation{Center for Quantitative Biology, Peking University, Beijing 100871, China}

\author{Yuhai Tu}
\email{ytu@flatironinstitute.org}
\affiliation{Center for Computational Biology \& Center for Computational Neuroscience, Flatiron Institute, New York, NY 10010, USA}

\date{\today}

\maketitle
\tableofcontents  
\clearpage
%% =========================================================
%% Formatting Tweaks for Supplement
%% =========================================================
% Reset counters and add "S" prefix
\setcounter{equation}{0}
\setcounter{figure}{0}
\setcounter{table}{0}
\setcounter{page}{1}
\setcounter{section}{0}

\renewcommand{\theequation}{S\arabic{equation}}
\renewcommand{\thefigure}{S\arabic{figure}}
\renewcommand{\thetable}{S\arabic{table}}

%% =========================================================
%% Content Starts Here
%% =========================================================

\section{Supplementary empirical experiments}
\subsection{Experimental setup}

We utilized a fixed subset of the MNIST dataset \cite{lecun1998gradient}, containing a training set of $N=1,000$ samples (100 randomly selected images per class) and a test set of $N=200$ samples. Input images were flattened into vectors of dimension $d_\mathrm{in}=784$. The model architecture consisted of a Multilayer Perceptron (MLP) with two hidden layers of width 50 (architecture: $784 \to 50 \to 50 \to 10$), employing ReLU activations. The model parameters were optimized using standard Stochastic Gradient Descent (SGD) without momentum or weight decay to minimize the cross-entropy loss.

To isolate the stochasticity introduced by batch sampling, we strictly enforced a deterministic initialization. Before the first run, weights were generated using the Kaiming Normal (He) initialization scheme and biases were set to zero. This configuration ($\bm{\theta}_0$) was saved and reloaded for all independent realizations. We conducted a systematic grid search over batch sizes $B \in \{10, 20, 50, 100, 200, 500, 1000\}$ (where $B=1000$ represents full-batch Gradient Descent) and learning rates $\eta \in \{0.001, 0.002, 0.005, 0.01, 0.02, 0.05, 0.1\}$. To ensure consistent convergence regimes, the total number of training iterations $T$ was scaled inversely with the learning rate according to $T = 100 / \eta$.

System trajectories were recorded at 101 evenly spaced time points $t \in [0, T]$. To analyze the landscape geometry efficiently, we tracked the specific weight matrix connecting the two hidden layers ($W^{(2)}$), treating it as a flattened parameter vector of dimension $d=2,500$. We computed the Hessian matrix $\mathbf{H}$ of the loss function restricted to this subspace at each checkpoint. Concurrently, we monitored functional performance (loss and accuracy) and tracked the indices of misclassified test samples to characterize sample-wise learning dynamics throughout the optimization process.

\subsection{Multi-valley landscape structure and Hessian spectrum}

Independent SGD runs initialized from the same deterministic starting point converge to multiple distinct low-loss endpoints that form separated clusters in parameter space. Fig.~\ref{fig:landscape_hessian}(a) visualizes this multi-valley structure by projecting the training trajectories for a representative subset of hyperparameter conditions onto the first two principal components of the collected endpoint weight vectors, with the vertical axis showing the training loss. Despite sharing the same initialization (star marker), trajectories under different noise conditions diverge and settle into distinct final clusters (open markers), confirming that the loss landscape contains multiple competing basins of attraction.

To quantify the functional dissimilarity between these endpoints, we computed the pairwise Jaccard similarity $\mathrm{Sim}_J$ (Eq.~3 of the main text) between the misclassified test sets of all converged endpoints across the full hyperparameter grid. The resulting similarity matrix, shown in Fig.~\ref{fig:landscape_hessian}(b), reveals a block-diagonal structure: endpoints trained under the same hyperparameter condition tend to be highly similar to one another, while endpoints from different conditions---particularly those differing in batch size---exhibit substantially lower similarity. This pattern indicates that the particular valley an SGD trajectory converges to is strongly influenced by the noise level during training.

We further verified that dissimilar endpoints are separated by loss barriers, confirming their assignment to distinct valleys. For each pair of converged endpoints $\bm{\theta}_A$ and $\bm{\theta}_B$, we evaluated the training loss along the linear interpolation path $\bm{\theta}(\alpha) = (1-\alpha)\bm{\theta}_A + \alpha\bm{\theta}_B$ for $\alpha \in [0,1]$ and recorded the barrier height $\Delta\mathcal{L}_\mathrm{barrier} = \max_\alpha \mathcal{L}(\bm{\theta}(\alpha)) - \max\{\mathcal{L}(\bm{\theta}_A), \mathcal{L}(\bm{\theta}_B)\}$. As illustrated in Fig.~\ref{fig:landscape_hessian}(c) and (d), endpoint pairs with high Jaccard similarity ($\mathrm{Sim}_J = 0.95$) share a flat interpolation path with negligible barriers, indicating they reside in the same valley. In contrast, pairs with low similarity ($\mathrm{Sim}_J = 0.625$) are separated by a prominent loss barrier, confirming they occupy distinct valleys. The negative correlation between the barrier height and the Jaccard similarity across all sampled pairs (Fig.~\ref{fig:landscape_hessian}(e); $R^2 = 0.64$) further validates the Jaccard similarity as a reliable proxy for basin identity.

To investigate the local geometry of these valleys, we analyzed the Hessian spectrum at the optimization endpoints. As shown in Fig.~\ref{fig:landscape_hessian}(f), the rank-ordered eigenvalue spectrum $\lambda_i(\mathbf{H})$ exhibits a rapid decay regardless of the batch size used during training. The vast majority of dimensions represent ``soft modes'' (near-zero curvature), implying that each solution lies on a manifold that is locally flat in most directions. To quantify this anisotropy, we calculated the fraction of the total curvature (trace of $\mathbf{H}$) captured by the top-10 eigenvalues for all converged realizations. The choice of $k=10$ is motivated by prior empirical analyses~\cite{sagunEmpiricalAnalysisHessian2018a}, which suggest that the Hessian spectrum of over-parameterized classification networks separates into a bulk of near-zero eigenvalues and a small set of large outliers roughly corresponding to the number of classes ($C=10$ for MNIST). As shown in Fig.~\ref{fig:landscape_hessian}(g), for most solutions these top 10 dimensions alone account for over 80\% of the total curvature, confirming that the optimization effectively confines the system to a low-dimensional active subspace governed by these principal directions.

\begin{figure*}[p]
\centering
\includegraphics[width=0.92\textwidth]{./Figures_SI/Fig1S_loss_landscape.pdf}
\caption{\textbf{Multi-valley landscape structure and Hessian spectrum analysis.}
(a)~PCA projection of SGD training trajectories for representative hyperparameter conditions. Trajectories share the same initialization (star) but diverge to distinct endpoint clusters (open markers); the vertical axis shows training loss.
(b)~Pairwise Jaccard similarity $\mathrm{Sim}_J$ between the misclassified test sets of all converged endpoints, grouped by hyperparameter condition ($\eta$, $B$).
(c,\,d)~Training loss along linear interpolation paths between endpoint pairs with high similarity ($\mathrm{Sim}_J = 0.95$, panel~c) and low similarity ($\mathrm{Sim}_J = 0.625$, panel~d). Line color indicates $t/t_\mathrm{max}$; $\Delta\mathcal{L}_\mathrm{barrier}$ denotes the barrier height.
(e)~Barrier height versus $\mathrm{Sim}_J$ for all sampled endpoint pairs ($R^2 = 0.64$).
(f)~Rank-ordered Hessian eigenvalue spectrum $\lambda_i(\mathbf{H})$ at endpoints, averaged over realizations for each batch size $B$.
(g)~Distribution of the fractional curvature captured by the top-10 eigenvalues across all converged solutions.}
\label{fig:landscape_hessian}
\end{figure*}

\subsection{Convergence stability and geometric properties of solutions}

To define the stability boundaries of the optimization, we mapped the convergence probability ($P_\mathrm{conv}$) across the hyperparameter grid (Fig.~\ref{fig:solution_stats}(a)). We identified a ``high-noise'' instability regime characterized by large learning rates and small batch sizes (e.g., $\eta \ge 0.05, B \le 50$), where the model consistently failed to achieve 100\% training accuracy. These configurations correspond to the excluded ``gray regions'' in the main text. Subsequent analyses are restricted to the stable regime.

We next examined the cosine similarity between solutions to further validate the diversity analysis. As shown in Fig.~\ref{fig:solution_stats}(b), low-noise settings produce highly collinear solution vectors ($\langle \mathrm{Sim}_{\cos} \rangle \approx 1$), whereas increasing the noise level leads to nearly orthogonal solutions. This observation is consistent with the Jaccard similarity results reported in the main text. This trend is accompanied by an increase in the weight displacement metric $\langle \|\Delta\bm{\theta}\|_2 \rangle$ (Fig.~\ref{fig:solution_stats}(c)), indicating that higher noise drives the final solutions significantly further from the initialization point $\bm{\theta}_0$.

Finally, we recorded the functional performance metrics for all configurations. As shown in Fig.~\ref{fig:solution_stats}(d) and (e), all converged runs achieve perfect training accuracy and near-zero training loss. However, the test loss landscape (Fig.~\ref{fig:solution_stats}(f)) exhibits significant variation across the grid, with moderate-to-high noise settings yielding lower test loss values compared to the full-batch baseline, mirroring the test accuracy trends discussed in the main text.

\begin{figure*}[htbp]
\centering
\includegraphics[width=1.0\textwidth]{./Figures_SI/Fig2S_solution_properties_supplement.pdf}
\caption{\textbf{Detailed statistics of convergence stability and solution geometry.} 
(a) Convergence probability ($P_\mathrm{conv}$), defined as the fraction of 20 independent runs that achieved 100\% training accuracy.
(b) Mean pairwise cosine similarity $\langle \mathrm{Sim}_{\cos} \rangle$ between final weight vectors.
(c) Mean Euclidean distance $\langle \|\Delta\bm{\theta}\|_2 \rangle$ of the final solution from the initialization $\bm{\theta}_0$.
(d--f) Heatmaps of mean final training accuracy (d), mean final training loss (e), and mean final test loss (f). Note that (e) and (f) are plotted on a logarithmic scale.}
\label{fig:solution_stats}
\end{figure*}

\subsection{Robustness of the freezing time definition}

As defined in the main text (Eq.~3), the freezing time $t_f$ is determined by the Jaccard similarity $\mathrm{Sim}_J$ between continuation endpoints: $t_f$ is the earliest checkpoint time after which all subsequent continuation endpoints remain mutually similar above a threshold $S_{\mathrm{th}}$. To verify that this definition is robust, we repeated the freezing time analysis of Fig.~2(a) using different similarity thresholds $S_{\mathrm{th}} \in \{0.90, 0.95, 0.99\}$. Fig.~\ref{fig:freeze_time_robustness} presents the corresponding heatmaps of the normalized mean freezing time $\eta \langle t_f \rangle$. Across all threshold choices, the qualitative dependence on hyperparameters remains unchanged: increasing SGD noise (larger $\eta$ or smaller $B$) systematically prolongs the transient phase. This robustness indicates that delayed freezing is an intrinsic dynamical feature of the optimization process, rather than an artifact of the particular similarity cutoff.

\begin{figure*}[htbp]
\centering
% TODO: replace with new figure generated using Jaccard-based freezing time
\includegraphics[width=1.0\textwidth]{./Figures_SI/Fig3S_freeze_time_supplement.pdf}
\caption{\textbf{Robustness of freezing time across different Jaccard similarity thresholds.}
Heatmaps of the normalized mean freezing time $\eta \langle t_f \rangle$ for different Jaccard similarity thresholds (a) $S_{\mathrm{th}} = 0.90$, (b) $S_{\mathrm{th}} = 0.95$, and (c) $S_{\mathrm{th}} = 0.99$.}
\label{fig:freeze_time_robustness}
\end{figure*}

\subsection{Details of continuation training}

To probe the properties of the basins explored by SGD, we adopted a deterministic continuation protocol. Starting from a standard SGD trajectory ($B=50, \eta=0.05$), we branched the training dynamics at selected time points $t_c$ and continued optimization using full-batch Gradient Descent ($B=1000$) for an additional 2,000 iterations. Because deterministic GD lacks stochastic noise, it cannot escape local minima; consequently, the final solution of each continuation branch effectively identifies the specific loss valley occupied by the optimizer at time $t_c$.

For the systematic grid scan reported in the main text (Fig.~2), we employed a matched continuation protocol: for each hyperparameter setting $(\eta, B)$, the deterministic continuation used the same learning rate $\eta$ as the SGD training trajectory, ensuring that the continuation dynamics probe the basin structure at the same effective gradient scale. This avoids artifacts from a mismatch between the training and probing learning rates. Continuation endpoints were compared using the Jaccard similarity of their misclassified test sets (Eq.~2 of the main text), and the freezing time $t_f$ was determined as defined in Eq.~3 of the main text.

As shown in Fig.~\ref{fig:continuation_details}(a--b) and (d--e), all continuation runs rapidly converge to near-zero training loss and achieve 100\% training accuracy. Despite this apparent convergence, the quality of the accessed basins varies substantially. Figures~\ref{fig:continuation_details}(c) and (f) summarize the corresponding generalization performance: solutions branched at early times ($t_c < 100$) exhibit poor test accuracy, whereas performance improves sharply throughout the transient phase and saturates around the freezing time $t_f$. This behavior indicates that SGD actively transitions from low-quality to high-quality basins before becoming dynamically trapped.

Importantly, this improvement in generalization is tightly linked to the geometry of the loss landscape. Figure~\ref{fig:continuation_details}(g) reveals a strong positive correlation between test accuracy and the flatness metric $F$ ($R^2 = 0.77$), providing further quantitative support for the conclusions drawn in Fig.~1(c) and~1(e) of the main text.

\begin{figure*}[htbp]
\centering
\includegraphics[width=1.0\textwidth]{./Figures_SI/Fig4S_escaping_behavior_supplement.pdf}
\caption{\textbf{Dynamics of continuation training.}
(a--c) Evolution of training loss (a), training accuracy (b), and test accuracy (c) for continuation runs branched at different times $t_c$. The solid blue curve denotes the reference SGD trajectory, while the colored curves correspond to deterministic GD continuations.
(d--f) Final training loss (d), training accuracy (e), and test accuracy (f) as functions of $t_c$.
(g) Correlation between solution flatness $F$ and final test accuracy, with $R^2 = 0.77$.}
\label{fig:continuation_details}
\end{figure*}

\subsection{Generalization to convolutional neural networks on CIFAR-10}

To demonstrate that our findings are not limited to simple fully connected networks or the MNIST dataset, we replicated the core experiments using a Convolutional Neural Network (CNN) trained on the CIFAR-10 dataset \cite{krizhevsky2009learning}.

\subsubsection{Experimental setup}
We utilized the CIFAR-10 dataset, which consists of 60,000 $32 \times 32$ color images in 10 classes. The model architecture, denoted as SimpleCNN, consists of three convolutional blocks followed by a classifier. Each convolutional block comprises a 2D convolution layer ($3 \times 3$ kernel, padding 1), a ReLU activation, and a $2 \times 2$ MaxPool layer. The channel depths for the three layers are 32, 64, and 64, respectively. The feature maps are flattened into a vector of dimension 1,024 and fed into a classifier consisting of a hidden fully connected layer with 128 units and a final output layer ($1024 \to 128 \to 10$).

We conducted a grid search over batch sizes $B \in \{20, 50, 100, 200, 500\}$ and learning rates $\eta \in \{0.002, 0.005, 0.01, 0.02, 0.05\}$. For each condition, 10 independent realizations were trained. To analyze the landscape geometry and solution diversity, we specifically focused on the weight matrix connecting the hidden layer to the output layer. This matrix corresponds to a parameter space of $d=1,280$ dimensions ($128 \times 10$).

\subsubsection{Results}
The results on CIFAR-10 corroborate the ``distinct valley'' hypothesis presented in the main text. Fig.~\ref{fig:cnn_cifar10}(a) visualizes the optimization dynamics in the projected weight space. Similar to the MNIST results, trajectories do not converge to a single global minimum but instead form isolated clusters, confirming that the loss landscape consists of multiple distinct valleys.

Notably, the separation between valleys is more pronounced for CIFAR-10 than for MNIST. As shown in Fig.~\ref{fig:cnn_cifar10}(b), the overall Jaccard similarity ($\mathrm{Sim}_J$) is substantially lower than that observed in the MNIST setting (Fig.~1(d) of the main text), suggesting that different local minima on the CIFAR-10 landscape correspond to functionally more distinct solutions, each failing on different subsets of hard examples.

Despite this increased landscape complexity, the relationship between noise and generalization remains consistent. As shown in Fig.~\ref{fig:cnn_cifar10}(c--f), increasing SGD noise (smaller $B$ or larger $\eta$) systematically drives the optimizer towards valleys that exhibit:
\begin{itemize}
    \item lower average Jaccard similarity, reflecting broader exploration across distinct valleys (Fig.~\ref{fig:cnn_cifar10}(c));
    \item higher test accuracy (Fig.~\ref{fig:cnn_cifar10}(d));
    \item increased solution flatness (Fig.~\ref{fig:cnn_cifar10}(e));
    \item a prolonged transient phase prior to convergence (Fig.~\ref{fig:cnn_cifar10}(f)).
\end{itemize}
Direct comparisons at a fixed learning rate ($\eta=0.005$) further confirm that reducing batch size effectively extends the transient training phase (Fig.~\ref{fig:cnn_cifar10}(g)) and enhances solution flatness (Fig.~\ref{fig:cnn_cifar10}(h)).

\begin{figure*}[hbtp]
\centering
\includegraphics[width=1.0\textwidth]{./Figures_SI/Fig5S_cnn_cifar10.pdf}
\caption{\textbf{Generalization of findings to a Convolutional Neural Network on CIFAR-10.} 
(a) A 2D Principal Component Analysis (PCA) projection of the weight training trajectories with different SGD hyperparameters (learning rate $\eta$ and batch size $B$), indicated by different colors and marker styles.
(b) The Jaccard similarity ($\mathrm{Sim}_J$) matrix between pairs of solutions.
(c--f) Heatmaps across the hyperparameter grid ($\eta, B$) for: 
(c) Average Jaccard similarity $\langle \mathrm{Sim}_J \rangle$ (lower values indicate greater diversity); 
(d) Maximum test accuracy $\max(Acc_{\mathrm{test}})$ ; 
(e) Maximum solution flatness $\max(F)$ ; 
(f) Normalized mean freezing time $\eta \langle t_f \rangle$. Gray regions in the heatmaps indicate divergent settings where no runs achieve 100\% training accuracy.
(g--h) Mean training loss $\langle \mathcal{L}_{\mathrm{train}} \rangle$ (g) and mean flatness $\langle F \rangle$ (h) across multiple runs for different batch sizes $B$ at a fixed learning rate $\eta=0.005$.}
\label{fig:cnn_cifar10}
\end{figure*}

\section{Supplementary information of the minimal two-valley model}
\subsection{Simulation settings}

The two-valley loss landscape $\mathcal{L}(x,y)$ is defined by Eq.~5 in the main text. The geometric properties of the landscape—specifically the bifurcation, valley depth, and flatness—are controlled by a set of fixed parameters. The values used in our simulations are listed in Table~\ref{tab:parameters}.

\begin{table}[h]
\centering
\caption{Parameters used for the minimal two-valley landscape model.}
\label{tab:parameters}
\renewcommand{\arraystretch}{1.2}
\begin{tabular}{ccl}
\hline
\textbf{Parameter} & \textbf{Value} & \textbf{Description} \\
\hline
$f_0$ & $1.0$ & Overall flatness scale factor \\
$x_0$ & $1.0$ & Normalization length scale \\
\hline
$x_1$ & $0.8$ & Position parameter for the flatter valley ($x > 0$) \\
$x_2$ & $0.65$ & Position parameter for the sharper valley ($x < 0$) \\
\hline
$y_f$ & $20$ & Scale controlling the rate of flatness growth \\
$y_b$ & $30$ & Scale controlling the valley bifurcation in $x$ \\
$y_d$ & $100$ & Decay scale for the gradient drift in $y$ \\
$\mathcal{L}_d$ & $1.0$ & Coefficient for the longitudinal drift speed \\
\hline
\end{tabular}
\end{table}

The SGD dynamics on this landscape were simulated using the discrete-time update rule given in Eq.~10 of the main text:
\begin{equation}
    \bm{\theta}_{t+1} = \bm{\theta}_t - \eta \left( \nabla \mathcal{L}(\bm{\theta}_t) + \bm{\xi}_t \right),
\end{equation}
where $\eta$ denotes the learning rate and $\bm{\theta}_t = (x_t, y_t)$ represents the parameters at iteration $t$. To capture the geometry-dependent nature of SGD noise, the stochastic term $\bm{\xi}_t$ was modeled as a zero-mean Gaussian vector with a covariance matrix proportional to the local curvature:
\begin{equation}
    \bm{\xi}_t \sim \mathcal{N}\bigl(0, \mathbf{\Sigma}(\bm{\theta}_t)\bigr), \quad \text{with} \quad \mathbf{\Sigma}(\bm{\theta}_t) = 2\sigma \mathbf{H}(\bm{\theta}_t).
\end{equation}
Here, $\sigma$ represents the noise strength. In empirical simulations, we employed the positive semi-definite rectification of the Hessian $\mathbf{H}$ by clipping negative eigenvalues to zero. This ensures numerical stability while preserving the anisotropic diffusion characteristics driven by the valley geometry.

We explored the parameter space by varying both $\eta$ and $\sigma$ across 10 logarithmically spaced values within the range $[10^{-2}, 10^{-1}]$. For each hyperparameter configuration $(\eta, \sigma)$, we performed $2000$ independent trials, each spanning $2\times10^4$ iterations. To mitigate initialization bias towards either valley, starting positions were fixed at $y_0 = 5$, with $x_0$ alternating between $-0.01$ and $+0.01$ across trials.

We quantified the dynamics using two key metrics:
\begin{enumerate}
    \item \textbf{Convergence probability ($P_\mathrm{flat}$):} The fraction of trajectories that converged to the flatter valley ($x > 0$) by the end of the simulation.
    \item \textbf{Freezing time ($t_f$):} The iteration index of the final sign change in the $x$-coordinate (i.e., the last crossing of the ridge at $x=0$).
\end{enumerate}
The resulting phase diagrams characterizing these dynamics are presented in Fig.~3(c) and~3(d) of the main text.

% \subsection{Results of the \texorpdfstring{$\Sigma\propto\mathbf{H}^2$}{Sigma proportional to H squared} case}

% Our analysis indicates that the selection of flat minima is driven by the anisotropic nature of SGD noise, where fluctuations are amplified in directions of high curvature. Consequently, the qualitative behavior of the system should not rely strictly on the linear proportionality $\mathbf{\Sigma} \propto \mathbf{H}$, provided that the noise strength remains monotonically increasing with curvature.

% To validate this robustness, we performed simulations using the quadratic scaling $\mathbf{\Sigma} \propto \mathbf{H}^2$, as derived from the ``Random Shift Model''~\cite{yangStochasticGradientDescent2023}. The results, presented in Fig.~, demonstrate that the convergence probability ($P_\mathrm{flat}$) and the freezing time ($t_f$) are qualitatively consistent with those obtained using the linear model. This confirms that our conclusions regarding stability and flatness selection are insensitive to the specific functional form of the geometry-dependent noise.

\subsection{Robustness under quadratic noise scaling}

Our analysis indicates that the selection of flat minima is driven by the anisotropic nature of SGD noise, where fluctuations are amplified in directions of high curvature. Consequently, the qualitative behavior of the system should not rely strictly on the linear proportionality $\mathbf{\Sigma} \propto \mathbf{H}$, provided that the noise strength remains monotonically increasing with curvature.

To validate this robustness, we performed simulations using the quadratic scaling $\mathbf{\Sigma} \propto \mathbf{H}^2$, as derived from the ``Random Shift Model''~\cite{yangStochasticGradientDescent2023}. The results, presented in Fig.~\ref{fig:toy_model_H2}, demonstrate that the convergence probability ($P_\mathrm{flat}$) and the freezing time ($t_f$) are qualitatively consistent with those obtained using the linear model. Specifically, increasing the learning rate $\eta$ or the noise strength $\sigma$ consistently biases the selection probability towards the flatter minimum (Fig.~\ref{fig:toy_model_H2}(a)) and prolongs the transient search phase (Fig.~\ref{fig:toy_model_H2}(b)). This confirms that our conclusions regarding stability and flatness selection are insensitive to the specific functional form of the geometry-dependent noise.

\begin{figure}[hbtp]
\centering
\includegraphics[width=1.0\linewidth]{./Figures_SI/Fig6S_toy_model_H2.pdf}
\caption{\textbf{Robustness of selection mechanism under quadratic noise scaling ($\mathbf{\Sigma} \propto \mathbf{H}^2$).}
(a) Probability $P_\mathrm{flat}$ of converging to the flat minimum as a function of learning rate $\eta$ and noise strength $\sigma$. 
(b) Normalized mean freezing time $\eta \langle t_f \rangle$.}
\label{fig:toy_model_H2}
\end{figure}

%% =========================================================
%% References for Supplement
%% =========================================================
% Create a separate .bib file for supplement or use the same one but separate the compilation.
% Usually, references solely cited in SI are listed here.
\clearpage
\bibliography{deeplearning}

\end{document}